\documentclass{article}

\usepackage[margin=1.5in]{geometry}
\usepackage{graphicx}
\usepackage{amsmath,amssymb,amsthm}
\usepackage{mathtools}
\usepackage{dsfont}
\usepackage[labelfont=bf]{caption}

\usepackage[dvipsnames,table]{xcolor}
\definecolor{c0}{HTML}{ffb000}
\definecolor{c1}{HTML}{fe6100}
\definecolor{c2}{HTML}{dc267f}
\definecolor{c3}{HTML}{785ef0}
\definecolor{c4}{HTML}{648fff}

\usepackage{hyperref}
\hypersetup{
    colorlinks=true,
    linkcolor=c2,
    citecolor=c4,
    filecolor=c4,
    urlcolor=c4
}
\usepackage[capitalize,nameinlink,noabbrev]{cleveref}
\crefformat{equation}{#2(#1)#3}
\crefmultiformat{equation}{(#2#1#3)}%
{ and~#2(#1)#3)}{, #2(#1)#3}{ and~#2(#1)#3}
\crefname{assumption}{Assumption}{Assumptions}
\crefname{proposition}{Proposition}{Propositions}

\usepackage{thmtools}
\usepackage{thm-restate}

\usepackage{algorithm}
\usepackage{algorithmicx}
\usepackage[noend]{algpseudocode}
\algrenewcommand\algorithmiccomment[1]{\hfill{\color{gray}$\triangleright$~#1}}

\usepackage[backend=biber,backref,style=alphabetic,maxalphanames=4,maxcitenames=99,maxbibnames=99,giveninits=true,url=false,doi=false,isbn=false,eprint=false,date=year]{biblatex}
\addbibresource{references.bib}
\DeclareSourcemap{
  \maps[datatype=bibtex, overwrite]{
    \map{
      \step[fieldset=editor, null]
    }
  }
}
\DefineBibliographyStrings{english}{%
  backrefpage = {cited on page},%
  backrefpages = {cited on pages}%
}

% ---------- macros (project conventions) ----------
\newcommand{\ones}{\mathds{1}}
\renewcommand{\vec}[1]{\mathbf{#1}}          % bold for vectors AND matrices
\newcommand{\T}{\mathsf{T}}                  % transpose: \vec{A}^\T
\newcommand{\F}{\mathrm{F}}                  % Frobenius subscript: \|\cdot\|_\F
\newcommand{\eps}{\varepsilon}
\newcommand{\cF}{\mathcal{F}}
\newcommand{\OPT}{\mathrm{OPT}}
\newcommand{\iopt}{i^\mathrm{opt}}
\newcommand{\inner}[2]{\langle #1, #2 \rangle}
\DeclareMathOperator{\tr}{tr}
\DeclareMathOperator{\rank}{rank}
\DeclareMathOperator{\orth}{orth}
\DeclareMathOperator*{\argmin}{arg\,min}

\newtheorem{theorem}{Theorem}[section]
\newtheorem{lemma}[theorem]{Lemma}
\newtheorem{corollary}[theorem]{Corollary}
\newtheorem{proposition}[theorem]{Proposition}
\newtheorem{assumption}[theorem]{Assumption}
\theoremstyle{remark}
\newtheorem{remark}[theorem]{Remark}

\title{Constant factor to (1+$\epsilon$)-approximation for structured matrix learning}
\author{}
\date{}

\newcommand{\COLT}[1]{Conference on Learning Theory (COLT) #1}


\begin{document}
\maketitle

\begin{abstract}
We study hypothesis selection in the matrix--vector query model.
A family $\cF = \{\vec{A}_1, \dots, \vec{A}_N\} \subset \mathbb{R}^{n \times n}$ is known, an unknown $\vec{A} \in \mathbb{R}^{n \times n}$ is accessible only through the products $\vec{x} \mapsto \vec{A}\vec{x}$ and $\vec{x} \mapsto \vec{A}^\T\vec{x}$, and the goal is to identify a member of $\cF$ that is nearly as close to $\vec{A}$ in Frobenius norm as the best one.
Estimating the $N$ distances separately costs $O(\eps^{-2}\log N)$ queries.
We show that, given a warm start --- any constant-factor approximation to $\vec{A}$, not necessarily drawn from $\cF$ --- only
\[
O\bigl( \eps^{-2}\sqrt{\log N} \bigr)
\]
queries are needed, a quadratic improvement in the dependence on $\log N$, with no dependence on the dimension $n$.

The saving comes from using both sides of the matrix.
Writing $k$ for the number of queries spent, a randomized rangefinder first deflates the residual between $\vec{A}$ and the warm start, leaving a remainder that is \emph{spectrally flat}: its spectral norm is a $1/\sqrt k$ fraction of its Frobenius norm.
A single shared Gaussian quadratic sketch then estimates all $N$ inner products against that remainder at once.
Flatness is what makes the sketch concentrate as though it drew $k^2$ samples rather than $k$, so a union bound over $N$ hypotheses is paid at $k = \Theta(\sqrt{\log N})$ instead of $\Theta(\log N)$.
Deflation is the step that requires transpose access, and we show this is unavoidable: in any bounded-precision model, algorithms restricted to one-sided queries need $\widetilde\Omega(\log N)$ queries even when handed a warm start.
\end{abstract}

\section{Model and main result}\label{sec:model}


Throughout, $\cF = \{\vec{A}_1, \dots, \vec{A}_N\} \subset \mathbb{R}^{n \times n}$ is known and $\vec{A} \in \mathbb{R}^{n \times n}$ is unknown.
The algorithm may issue \emph{right queries} $\vec{x} \mapsto \vec{A}\vec{x}$ and \emph{left queries} $\vec{x} \mapsto \vec{A}^\T\vec{x}$ with arbitrary real query vectors $\vec{x} \in \mathbb{R}^n$, and observes the full output vector exactly.
Each query vector may depend on the answers to all previous queries.
The query complexity is the total number of matrix--vector products issued; all other computation (in particular, on the known family $\cF$) is free.
We write
\begin{equation}\label{eq:opt}
d_i \coloneqq \|\vec{A} - \vec{A}_i\|_\F,
\qquad
\OPT \coloneqq \min_{i \in [N]} d_i,
\qquad
\iopt \in \argmin_{i \in [N]} d_i .
\end{equation}
The goal is to output an index $i^\star \in [N]$ with $d_{i^\star} \le (1+\eps)\,\OPT$.

We use $\|\cdot\|_2$ for the spectral norm, $\|\cdot\|_*$ for the nuclear norm, $\sigma_1(\vec{X}) \ge \sigma_2(\vec{X}) \ge \cdots$ for singular values, and $[\vec{X}]_r$ for a best rank-$r$ approximation of $\vec{X}$ in Frobenius norm.
The inner product is $\inner{\vec{X}}{\vec{Y}} = \tr(\vec{X}^\T\vec{Y})$.

We assume throughout that a \emph{warm start} is available.

\begin{assumption}[Warm start]\label{asm:warm}
A matrix $\bar{\vec{A}}_0 \in \mathbb{R}^{n \times n}$ with known entries, not necessarily in $\cF$, is given, together with a constant $\gamma \ge 1$ such that $\|\vec{A} - \bar{\vec{A}}_0\|_\F \le \gamma\,\OPT$.
\end{assumption}

The constant $\gamma$ need not be tight; only the bound is used.
Note that \cref{asm:warm} is nontrivial exactly when $\OPT > 0$: if $\OPT = 0$ then $\bar{\vec{A}}_0 = \vec{A}$ and the problem is solved offline.

\begin{theorem}[Main]\label{thm:main}
Let $\eps \in (0,1]$ and $\delta \in (0,1/2)$, and let \cref{asm:warm} hold.
There is an algorithm making at most
\[
O\!\left( \gamma^4\eps^{-2}\sqrt{\log(N/\delta)} \;+\; \log(1/\delta) \right)
\]
queries that outputs $i^\star$ with $d_{i^\star} \le (1+\eps)\OPT$ with probability at least $1-\delta$.
The query count is independent of $n$.
\end{theorem}

For constant $\gamma$ and $\delta$ this reads $O(\eps^{-2}\sqrt{\log N})$, the form quoted in the abstract.
The algorithm uses only two rounds of queries: all right queries are made first, with vectors drawn independently of $\vec{A}$, and the left queries follow, with vectors spanning the range of the first round's output.

The argument has three moving parts.
\Cref{sec:reduction} reduces selection to estimating $N$ inner products against the residual, and shows how accurate those estimates must be.
\Cref{sec:estimate} gives the estimator and its analysis; this is where the $\sqrt{\log N}$ appears.
\Cref{sec:mainproof} combines the two.
\Cref{sec:discussion} explains why two-sided access is necessary (\cref{sec:onesided}) and why deflation is (\cref{sec:whydeflation}).

\section{Preliminaries}\label{sec:prelim}

We need two standard tools: Hanson--Wright concentration for Gaussian quadratic forms, and the spectral-norm guarantee of the randomized rangefinder.

\begin{lemma}[Hanson--Wright, averaged form \cite{Hanson1971,rudelson2013hanson}]\label{lem:hw}
Let $\vec{N} \in \mathbb{R}^{n \times n}$ be fixed and let $\vec{y}_1, \dots, \vec{y}_k \sim \mathcal{N}(\vec 0, \vec{I}_n)$ be independent.
Then for every $t \ge 0$,
\begin{equation}\label{eq:hw}
\Pr\left[ \left| \frac{1}{k}\sum_{b=1}^{k} \vec{y}_b^\T \vec{N} \vec{y}_b - \tr \vec{N} \right| \ge t \right]
\le 2\exp\left( -\Omega\!\left( \min\left\{ \frac{k t^2}{\|\vec{N}\|_\F^2},\; \frac{k t}{\|\vec{N}\|_2} \right\} \right) \right).
\end{equation}
\end{lemma}


We emphasize that \cref{lem:hw} does not require $\vec N$ to be symmetric or positive semidefinite.
Indeed, $\vec{y}^\T\vec{N}\vec{y} = \vec{y}^\T \vec{N}_{\mathrm{sym}} \vec{y}$ and $\tr\vec{N} = \tr\vec{N}_{\mathrm{sym}}$ for the symmetric part $\vec{N}_{\mathrm{sym}} \coloneqq \tfrac12(\vec{N} + \vec{N}^\T)$, so the symmetric case applies verbatim; and since $\|\vec{N}_{\mathrm{sym}}\|_\F \le \|\vec{N}\|_\F$ and $\|\vec{N}_{\mathrm{sym}}\|_2 \le \|\vec{N}\|_2$ by the triangle inequality, replacing $\vec{N}_{\mathrm{sym}}$ by $\vec{N}$ in the exponent of \cref{eq:hw} only weakens the bound.
This is used with $\vec{N}_i = (\vec{B}_i^H)^\T\vec{R}$ in \cref{sec:estimate}, which is not symmetric.

\begin{lemma}[Deflation via the randomized rangefinder \cite{HMT11}]\label{lem:rsvd}
Let $\vec{M} \in \mathbb{R}^{n \times n}$, let $\vec{\Omega} \in \mathbb{R}^{n \times k}$ have i.i.d.\ $\mathcal{N}(0,1)$ entries with $k \ge 4$, and let $\vec{Q} = \orth(\vec{M}\vec{\Omega})$.
Then with probability $1 - e^{-\Omega(k)}$,
\begin{equation}\label{eq:deflation}
\|(\vec{I} - \vec{Q}\vec{Q}^\T)\vec{M}\|_2 = O\!\left( \frac{\|\vec{M}\|_\F}{\sqrt{k}} \right) .
\end{equation}
\end{lemma}

\noindent
In words: sketching with $k$ Gaussian vectors leaves a residual whose spectral norm is a $1/\sqrt k$ fraction of the Frobenius norm of $\vec M$ --- the residual is \emph{spectrally flat}.
This is the only property of the rangefinder we use.

\begin{proof}
The claim is trivial for $k = O(1)$, since $\|(\vec{I} - \vec{Q}\vec{Q}^\T)\vec{M}\|_2 \le \|\vec{M}\|_\F$ always; so assume $k$ exceeds any fixed constant.
Write $\sigma_j \coloneqq \sigma_j(\vec M)$ and split $k$ into a target rank $\rho \coloneqq \lfloor k/2\rfloor - 1$ and an oversampling parameter $p \coloneqq k - \rho$, so that $\rho \le p$ and $p = \Theta(k)$.
The deviation bound for the Gaussian rangefinder \cite[Thm.~10.8]{HMT11}, applied with $t = e$ and $u = \sqrt k$, gives with probability $1 - e^{-\Omega(k)}$
\[
\|(\vec{I} - \vec{Q}\vec{Q}^\T)\vec{M}\|_2
\le \Bigl( 1 + e\sqrt{12\rho/p} + \frac{e^2 \sqrt{k}\,\sqrt{\rho+p}}{p+1} \Bigr) \sigma_{\rho+1}
+ \frac{e^2 \sqrt{\rho + p}}{p+1} \Bigl( \sum_{j > \rho} \sigma_j^2 \Bigr)^{1/2}.
\]
Since $\rho \le p$ and $\rho + p = k = O(p)$, the coefficient of $\sigma_{\rho+1}$ is $O(1)$ and that of the tail is $O(1/\sqrt k)$.
The tail is at most $\|\vec M\|_\F$, while $\rho + 1 = \Theta(k)$ forces $\sigma_{\rho+1} \le \|\vec{M}\|_\F/\sqrt{\rho+1} = O(\|\vec M\|_\F/\sqrt k)$.
Both terms are therefore $O(\|\vec M\|_\F/\sqrt k)$.
\end{proof}

We record two elementary facts used repeatedly.
For any $\vec{X}$ with $\rank(\vec{X}) \le r$ we have $\|\vec{X}\|_* \le \sqrt{r}\,\|\vec{X}\|_\F$.
For any $\vec{X}$ and any $r \ge 0$, the tail $\vec{X} - [\vec{X}]_r$ satisfies $\|\vec{X} - [\vec{X}]_r\|_2 = \sigma_{r+1}(\vec{X}) \le \|\vec{X}\|_\F / \sqrt{r+1}$.

\section{Reduction: selection from approximate linear functionals}\label{sec:reduction}

Fix a \emph{pivot} $\bar{\vec{A}} \in \mathbb{R}^{n \times n}$ whose entries are known to the algorithm (it need not lie in $\cF$), and write
\[
\vec{B} \coloneqq \vec{A} - \bar{\vec{A}},
\qquad
\vec{B}_i \coloneqq \vec{A}_i - \bar{\vec{A}} \quad (i \in [N]).
\]
Expanding $d_i^2 = \|\vec{B} - \vec{B}_i\|_\F^2$ gives
\begin{equation}\label{eq:decomp}
d_i^2 = \|\vec{B}\|_\F^2 - 2\inner{\vec{B}}{\vec{B}_i} + \|\vec{B}_i\|_\F^2 .
\end{equation}
The term $\|\vec{B}\|_\F^2$ is unknown but common to all $i$, and $\|\vec{B}_i\|_\F^2$ is computable offline.
So selection reduces to estimating $\inner{\vec{B}}{\vec{B}_i}$ for all $i \in [N]$.
The following lemma guarantees that, for a sufficiently accurate pivot and with sufficiently accurate estimates, we can easily obtain a near-optimal approximation.


\begin{lemma}[Selection from relatively accurate functionals]\label{lem:select}
Let $\eps \in (0,1]$ and $\gamma \ge 1$, and let $\bar{\vec{A}}$ be a pivot with $\|\vec{B}\|_\F = \|\vec{A} - \bar{\vec{A}}\|_\F \le \gamma\,\OPT$.
Set
\begin{equation}\label{eq:eps0}
\eps_0 \coloneqq \frac{\eps}{2\gamma(1+\gamma)} \in (0,1] .
\end{equation}
Suppose $\hat{\varphi}_1, \dots, \hat{\varphi}_N$ satisfy
\begin{equation}\label{eq:relacc}
\bigl|\hat{\varphi}_i - \inner{\vec{B}}{\vec{B}_i}\bigr| \le \eps_0 \, \|\vec{B}_i\|_\F \, \|\vec{B}\|_\F
\qquad\text{for every } i \in [N],
\end{equation}
and let $i^\star \in \argmin_i \bigl( \|\vec{B}_i\|_\F^2 - 2\hat{\varphi}_i \bigr)$.
Then $d_{i^\star} \le \bigl(1 + \eps\bigr)\,\OPT$.
\end{lemma}

\begin{proof}
Write $\eta_i \coloneqq \eps_0\|\vec{B}_i\|_\F\|\vec{B}\|_\F$ for the error bound in \cref{eq:relacc}, and $s_i \coloneqq \|\vec{B}_i\|_\F^2 - 2\hat{\varphi}_i$ for the score minimized by $i^\star$.

\emph{Step 1: selection.}
By \cref{eq:decomp}, every $i \in [N]$ satisfies the identity $d_i^2 = \|\vec{B}\|_\F^2 + s_i + 2\bigl(\hat{\varphi}_i - \inner{\vec{B}}{\vec{B}_i}\bigr)$.
Applying it at $i = i^\star$ and using \cref{eq:relacc}, then the minimality $s_{i^\star} \le s_{\iopt}$, then the identity again at $i = \iopt$ and \cref{eq:relacc} once more,
\begin{equation}\label{eq:selection}
d_{i^\star}^2
\le \|\vec{B}\|_\F^2 + s_{i^\star} + 2\eta_{i^\star}
\le \|\vec{B}\|_\F^2 + s_{\iopt} + 2\eta_{i^\star}
\le d_{\iopt}^2 + 2\eta_{\iopt} + 2\eta_{i^\star}
= \OPT^2 + 2\eta_{\iopt} + 2\eta_{i^\star} .
\end{equation}
The unknown quantity $\|\vec{B}\|_\F^2$ cancels; only differences of the scores $s_i$ are used.

If $\|\vec{B}\|_\F = 0$ then $\eta_i = 0$ for every $i$ and \cref{eq:selection} gives $d_{i^\star} \le \OPT$.
Assume henceforth $\|\vec{B}\|_\F > 0$, so also $\OPT > 0$.

\emph{Step 2: the quadratic.}
The triangle inequality gives $\|\vec{B}_i\|_\F \le d_i + \|\vec{B}\|_\F$, so \cref{eq:selection} becomes
\[
d_{i^\star}^2
\le \OPT^2 + 2\eps_0 \|\vec{B}\|_\F \bigl( \|\vec{B}_{i^\star}\|_\F + \|\vec{B}_{\iopt}\|_\F \bigr)
\le \OPT^2 + 2\eps_0\|\vec{B}\|_\F\,\bigl( d_{i^\star} + \OPT + 2\|\vec{B}\|_\F \bigr).
\]
Write $a \coloneqq 2\eps_0\|\vec{B}\|_\F > 0$, so that $x \coloneqq d_{i^\star}$ obeys $x^2 - a x - \bigl( \OPT^2 + a\OPT + 2a\|\vec{B}\|_\F \bigr) \le 0$.
Since the leading coefficient is positive, $x$ is at most the larger root of the associated quadratic, namely
\[
x \;\le\; \frac{a}{2} + \sqrt{ \frac{a^2}{4} + \OPT^2 + a\OPT + 2a\|\vec{B}\|_\F } 
= \frac{a}{2} + \sqrt{u^2+v},
\]
where we have introduced the notation $u \coloneqq \OPT + \tfrac{a}{2} > 0$ and $v \coloneqq 2 a \|\vec{B}\|_\F > 0$.
Now $\bigl( u + \tfrac{v}{2u} \bigr)^2 = u^2 + v + \tfrac{v^2}{4u^2} \ge u^2 + v$, and hence $\sqrt{u^2 + v} \le u + \tfrac{v}{2u}$.
Using this, we obtain a bound
\[
x \;\le\; \frac{a}{2} + u + \frac{v}{2u}
\;=\; \frac{a}{2} + \Bigl( \OPT + \frac{a}{2} \Bigr) + \frac{a\|\vec{B}\|_\F}{\OPT + a/2}
\;=\; \OPT + a + \frac{a\|\vec{B}\|_\F}{\OPT + a/2} .
\]
Discarding the term $a/2 \ge 0$ from the last denominator gives
\begin{equation}\label{eq:pivotfree}
d_{i^\star} \;\le\; \OPT + a\Bigl( 1 + \frac{\|\vec{B}\|_\F}{\OPT} \Bigr)
= \OPT + 2\eps_0 \|\vec{B}\|_\F \Bigl( 1 + \frac{\|\vec{B}\|_\F}{\OPT} \Bigr).
\end{equation}

\emph{Step 3: the warm start.}
By assumption $\|\vec{B}\|_\F \le \gamma\OPT$ so by our choice of $\eps_0 = \eps / (2\gamma(1+\gamma))$,
\[
d_{i^\star} \;\le\; \OPT + 2\eps_0\gamma(1+\gamma)\,\OPT \;=\; \bigl(1 + \eps \bigr)\OPT . \qedhere
\]
\end{proof}


\section{The estimation routine}\label{sec:estimate}

The routine below is the workhorse, invoked with a pivot $\bar{\vec{A}}$, an accuracy $\eps_0 \in (0,1]$, and a failure probability $\delta$.
Its structure is that of a Hutch\texttt{++}-style estimator: deflate a low-rank part, evaluate its contribution exactly, and sketch only the flat remainder.
It differs in where the deflation is applied.
A direct transcription would deflate each product $\vec{B}_i^\T\vec{B}$, whose trace is the target $\inner{\vec{B}}{\vec{B}_i}$; that would require a separate rangefinder per candidate, and hence $\Omega(N)$ queries.
Instead we deflate the single unknown factor $\vec{B}$ once, at a cost of $O(k)$ queries shared across all $N$ candidates, and flatten the other factor $\vec{B}_i$ offline by truncation --- which is free, because $\vec{B}_i$ is known.
The product $\vec{N}_i = (\vec{B}_i^H)^\T\vec{R}$ is then flat because both of its factors are, which is what the union bound over $N$ candidates needs.

\medskip
\noindent\textbf{Algorithm \textsc{EstimateScores}$(\bar{\vec{A}}, \eps_0, \delta)$.}
Set
\begin{equation}\label{eq:kdef}
k \coloneqq \Bigl\lceil C\bigl( \eps_0^{-2}\sqrt{\log(N/\delta)} + \log(1/\delta) \bigr) \Bigr\rceil,
\qquad
r_0 \coloneqq \bigl\lceil \alpha\, \eps_0^{2} k \bigr\rceil,
\end{equation}
where $C$ is a sufficiently large absolute constant and $\alpha \in (0,1)$ is a sufficiently small one.
Thus $k = O\bigl( \eps_0^{-2}\sqrt{\log(N/\delta)} + \log(1/\delta) \bigr)$ and $r_0 = \Theta(\eps_0^2 k)$.
The truncation rank $r_0$ trades the two error terms against each other: shrinking $\alpha$ improves the low-rank bound \cref{eq:lowrank} and worsens the flat bound, and the analysis fixes $\alpha$ first, then takes $C$ large.
\begin{enumerate}
\item \textbf{(Sketch.)}
Draw $\vec{\Omega}, \vec{Y} \in \mathbb{R}^{n \times k}$ with i.i.d.\ $\mathcal{N}(0,1)$ entries, independent of each other, and issue the $2k$ right queries $\vec{A}\vec{\Omega}$ and $\vec{A}\vec{Y}$.
Form $\vec{B}\vec{\Omega} = \vec{A}\vec{\Omega} - \bar{\vec{A}}\vec{\Omega}$ and $\vec{Q} = \orth(\vec{B}\vec{\Omega}) \in \mathbb{R}^{n \times k'}$, $k' \le k$.
\item \textbf{(Deflate.)}
Issue the $k'$ left queries $\vec{A}^\T\vec{q}_j$, $j \in [k']$, and form $\vec{W} \coloneqq \vec{Q}^\T\vec{B} = (\vec{A}^\T\vec{Q})^\T - \vec{Q}^\T\bar{\vec{A}}$ and $\widehat{\vec{B}} \coloneqq \vec{Q}\vec{W} = \vec Q\vec Q^\T\vec B$.
Write $\vec{R} \coloneqq \vec{B} - \widehat{\vec{B}} = (\vec{I} - \vec{Q}\vec{Q}^\T)\vec{B}$ (unknown to the algorithm).
\item \textbf{(Split the candidates.)}
For each $i$, compute the SVD of $\vec{B}_i$ and split
\[
\vec{B}_i = \vec{B}_i^L + \vec{B}_i^H,
\qquad
\vec{B}_i^L \coloneqq [\vec{B}_i]_{r_0},
\]
so that $\rank(\vec{B}_i^L) \le r_0$ and $\|\vec{B}_i^H\|_2 = \sigma_{r_0+1}(\vec B_i) \le \|\vec{B}_i\|_\F / \sqrt{r_0 + 1}$.
\item \textbf{(Quadratic sketch.)}
Using the queries $\vec{A}\vec{Y}$ already made in step 1, form $\vec{R}\vec{Y} =(\vec{A}\vec{Y} - \bar{\vec{A}}\vec{Y}) - \widehat{\vec{B}}\vec{Y}$ and, for each $i$,
\[
\hat{\psi}_i \coloneqq \frac{1}{k} \inner{\vec{B}_i^H \vec{Y}}{\vec{R}\vec{Y}} = \frac{1}{k}\sum_{b=1}^{k} \vec{y}_b^\T (\vec{B}_i^H)^\T \vec{R}\, \vec{y}_b .
\]
\item \textbf{(Output.)}
Return $\hat{\varphi}_i \coloneqq \inner{\vec{B}_i}{\widehat{\vec{B}}} + \hat{\psi}_i$ for all $i \in [N]$.
\end{enumerate}
The query cost is $2k$ right queries and $k' \le k$ left queries, i.e.\ $O(k)$ in total, issued in two rounds: the right queries of step 1 use vectors drawn independently of $\vec{A}$, and only the left queries of step 2 depend on answers already received.
Everything else --- $\bar{\vec{A}}\vec{\Omega}$, $\vec{Q}^\T\bar{\vec{A}}$, $\bar{\vec{A}}\vec{Y}$, the SVDs of the $\vec{B}_i$, and $\vec{B}_i^H\vec{Y}$ --- is computed offline.
In particular $\vec R\vec Y$ is computable even though $\vec R$ is not.

\begin{lemma}[Guarantee for \textsc{EstimateScores}]\label{lem:estimate}
For a sufficiently large absolute constant $C$ in \cref{eq:kdef}, the following holds.
For any pivot $\bar{\vec{A}}$, any $\eps_0 \in (0,1]$, and any $\delta \in (0,1)$, with probability at least $1 - \delta$ the output of \textsc{EstimateScores} satisfies
\[
\bigl| \hat{\varphi}_i - \inner{\vec{B}}{\vec{B}_i} \bigr| \le \eps_0 \, \|\vec{B}_i\|_\F \, \|\vec{B}\|_\F
\qquad\text{simultaneously for every } i \in [N].
\]
\end{lemma}

\begin{proof}
\emph{Degenerate cases.}
If $\|\vec{B}\|_\F = 0$ then $\vec{B}\vec{\Omega} = \vec 0$, so $k' = 0$, $\widehat{\vec{B}} = \vec 0$, $\vec R = \vec 0$, $\hat\psi_i = 0$, and $\hat\varphi_i = 0 = \inner{\vec{B}}{\vec{B}_i}$ for all $i$; the claim holds with probability $1$.
Similarly, for any index $i$ with $\|\vec{B}_i\|_\F = 0$ we have $\vec{B}_i = \vec 0$, hence $\vec{B}_i^L = \vec{B}_i^H = \vec 0$, hence $\hat\varphi_i = 0 = \inner{\vec{B}}{\vec{B}_i}$; such indices may be ignored below.
So assume $\|\vec{B}\|_\F > 0$ and restrict attention to $i$ with $\|\vec{B}_i\|_\F > 0$.

\emph{Deterministic facts.}
Because $\vec{Q}$ has orthonormal columns, $\vec{I} - \vec{Q}\vec{Q}^\T$ is an orthogonal projector, so $\vec R = (\vec I - \vec Q\vec Q^\T)\vec B$ satisfies
\begin{equation}\label{eq:orth}
\|\vec{R}\|_\F \le \|\vec{B}\|_\F
\end{equation}
for every realization of $\vec{\Omega}$.
Also, $\|\vec{B}_i^L\|_\F \le \|\vec{B}_i\|_\F$ and $\|\vec{B}_i^H\|_\F \le \|\vec{B}_i\|_\F$.

\emph{Error decomposition.}
Since $\vec{B} = \widehat{\vec{B}} + \vec{R}$ and $\vec{B}_i = \vec{B}_i^L + \vec{B}_i^H$,
\[
\inner{\vec{B}}{\vec{B}_i} - \inner{\vec{B}_i}{\widehat{\vec{B}}}
= \inner{\vec{B}_i}{\vec{R}}
= \inner{\vec{B}_i^L}{\vec{R}} + \inner{\vec{B}_i^H}{\vec{R}},
\]
and therefore, deterministically,
\begin{equation}\label{eq:errsplit}
\bigl| \hat{\varphi}_i - \inner{\vec{B}}{\vec{B}_i} \bigr|
= \bigl|  \bigl( \hat{\psi}_i - \inner{\vec{B}_i^H}{\vec{R}} \bigr) - \inner{\vec{B}_i^L}{\vec{R}} \bigr|
\leq \bigl|  \hat{\psi}_i - \inner{\vec{B}_i^H}{\vec{R}} \bigr| + \bigl| \inner{\vec{B}_i^L}{\vec{R}} \bigr|.
\end{equation}
The two brackets are bounded separately.

\emph{Deflation event.}
Let $E_1$ be the event that \cref{eq:deflation} holds, i.e.\ $\|\vec{R}\|_2 = O\bigl( \|\vec{B}\|_\F/\sqrt{k} \bigr)$.
Since $k = \Omega(\log(1/\delta))$ by \cref{eq:kdef}, \cref{lem:rsvd} gives $\Pr[E_1^c] \le e^{-\Omega(k)} \le \delta/4$ for $C$ large enough.

\emph{Low-rank part: deterministic given $E_1$.}
Since $\rank(\vec{B}_i^L) \le r_0$,
\[
|\inner{\vec{B}_i^L}{\vec{R}}|
\le \|\vec{B}_i^L\|_* \, \|\vec{R}\|_2
= O\!\left( \sqrt{\frac{r_0}{k}} \right) \|\vec{B}_i\|_\F \, \|\vec{B}\|_\F
\qquad\text{on } E_1 .
\]
By \cref{eq:kdef}, $\sqrt{r_0/k} = O(\sqrt{\alpha}\,\eps_0)$, so choosing $\alpha$ a small enough absolute constant gives
\begin{equation}\label{eq:lowrank}
|\inner{\vec{B}_i^L}{\vec{R}}| \le \frac{\eps_0}{4} \, \|\vec{B}_i\|_\F \, \|\vec{B}\|_\F
\qquad\text{on } E_1, \text{ for every } i .
\end{equation}

\emph{Flat part: Hanson--Wright.}
Condition on $\vec{\Omega}$; then $\vec{Q}$, $\vec{R}$, and every $\vec{B}_i^H$ are fixed, and $\vec{Y}$ is an independent Gaussian matrix whose columns $\vec{y}_1,\dots,\vec{y}_k$ remain i.i.d.\ $\mathcal{N}(\vec 0,\vec{I}_n)$ after conditioning.
So \cref{lem:hw} applies to each $i$ with the fixed matrix $\vec{N}_i \coloneqq (\vec{B}_i^H)^\T \vec{R}$: indeed
\[
\hat{\psi}_i = \frac{1}{k}\sum_{b=1}^{k} \vec{y}_b^\T \vec{N}_i \vec{y}_b ,
\qquad
\tr(\vec{N}_i) = \inner{\vec{B}_i^H}{\vec{R}} ,
\]
so $\hat\psi_i$ is an averaged quadratic form whose mean is exactly the quantity it must estimate in \cref{eq:errsplit}.
(Here $\vec{N}_i$ is not symmetric, which \cref{lem:hw} permits.)
Because $r_0 + 1 \ge \alpha\eps_0^2 k$ we have $\|\vec{B}_i^H\|_2 \le \|\vec{B}_i\|_\F/\sqrt{r_0+1} = O\bigl( \|\vec{B}_i\|_\F/(\eps_0\sqrt{k}) \bigr)$, the implied constant depending only on the now-fixed $\alpha$; so on $E_1$,
\[
\|\vec{N}_i\|_\F \le \|\vec{B}_i^H\|_2 \|\vec{R}\|_\F = O\!\left( \frac{\|\vec{B}_i\|_\F \|\vec{B}\|_\F}{\eps_0 \sqrt{k}} \right),
\qquad
\|\vec{N}_i\|_2 \le \|\vec{B}_i^H\|_2 \|\vec{R}\|_2 = O\!\left( \frac{\|\vec{B}_i\|_\F \|\vec{B}\|_\F}{\eps_0 k} \right),
\]
using $\|\vec{R}\|_\F \le \|\vec{B}\|_\F$ from \cref{eq:orth} for the first bound.
Apply \cref{lem:hw} with $t_i \coloneqq \tfrac{\eps_0}{2}\|\vec{B}_i\|_\F \|\vec{B}\|_\F > 0$, and abbreviate $\Pi \coloneqq \|\vec{B}_i\|_\F\|\vec{B}\|_\F$, so that $t_i = \tfrac{\eps_0}{2}\Pi$; the two exponents satisfy
\[
\frac{k t_i^2}{\|\vec{N}_i\|_\F^2}
= \Omega\!\left( \frac{k \cdot \eps_0^2\Pi^2}{\Pi^2/(\eps_0^2 k)} \right)
= \Omega\bigl( \eps_0^4 k^2 \bigr),
\qquad\quad
\frac{k t_i}{\|\vec{N}_i\|_2}
= \Omega\!\left( \frac{k \cdot \eps_0\Pi}{\Pi/(\eps_0 k)} \right)
= \Omega\bigl( \eps_0^2 k^2 \bigr) = \Omega\bigl( \eps_0^4 k^2 \bigr),
\]
the last step because $\eps_0 \le 1$; note that the scale $\Pi$ cancels, so one bound serves all $N$ candidates.
Thus the minimum in \cref{eq:hw} is $\Omega(\eps_0^4k^2)$, and a union bound over the $N$ candidates is affordable exactly when this absorbs $\log(N/\delta)$, i.e.\ when
\[
\eps_0^4 k^2 = \Omega\bigl( \log(N/\delta) \bigr)
\qquad\Longleftrightarrow\qquad
k = \Omega\bigl( \eps_0^{-2}\sqrt{\log(N/\delta)} \bigr) ,
\]
which is how $k$ was chosen in \cref{eq:kdef}: the exponent is quadratic in $k$, so solving for $k$ takes a square root of $\log N$.
Concretely, $\eps_0^2 k \ge C\sqrt{\log(N/\delta)}$ by \cref{eq:kdef} makes both exponents $\Omega\bigl( C^2\log(N/\delta) \bigr)$.
With $\alpha$ already fixed, taking $C$ large enough makes the right side of \cref{eq:hw} at most $\delta/(4N)$.
Thus, for every realization of $\vec{\Omega}$ in $E_1$, the conditional probability of $|\hat{\psi}_i - \inner{\vec{B}_i^H}{\vec{R}}| > t_i$ is at most $\delta/(4N)$, and a union bound over $i \in [N]$ gives
\[
\Pr\Bigl[ E_1 \cap \bigcup_{i \in [N]} \bigl\{ |\hat{\psi}_i - \inner{\vec{B}_i^H}{\vec{R}}| > t_i \bigr\} \Bigr] \le \frac{\delta}{4} .
\]

\emph{Conclusion.}
With probability at least $1 - \Pr[E_1^c] - \delta/4 \ge 1 - \delta/2 \ge 1 - \delta$, both \cref{eq:lowrank} and $|\hat{\psi}_i - \inner{\vec{B}_i^H}{\vec{R}}| \le t_i$ hold for all $i$, and then \cref{eq:errsplit} gives
\[
\bigl|\hat{\varphi}_i - \inner{\vec{B}}{\vec{B}_i}\bigr| \le \frac{\eps_0}{2}\|\vec{B}_i\|_\F\|\vec{B}\|_\F + \frac{\eps_0}{4}\|\vec{B}_i\|_\F\|\vec{B}\|_\F \le \eps_0 \|\vec{B}_i\|_\F\|\vec{B}\|_\F . \qedhere
\]
\end{proof}

\begin{remark}[Where $\sqrt{\log N}$ comes from]\label{rem:k2}
The union bound in the proof is paid against exponents of order $\eps_0^4 k^2$, i.e.\ against an \emph{effective sample size of $k^2$} from only $O(k)$ queries.
Both factors of $k$ have distinct origins: one from the averaging over the $k$ columns of $\vec{Y}$, and one from the spectral flatness of $\vec{N}_i = (\vec{B}_i^H)^\T\vec{R}$, whose two factors were flattened separately --- $\vec{B}_i^H$ by the rank-$r_0$ truncation (legal because $\vec{B}_i$ is known), and $\vec{R}$ by the rangefinder (legal because left queries make $\vec{Q}^\T\vec{B}$ exactly computable).
Setting $\eps_0^4 k^2 = \Theta(\log N)$ gives $k = \Theta\bigl(\eps_0^{-2}\sqrt{\log N}\bigr)$.
Without the two flattening steps the exponent degrades to order $\eps_0^2 k$, forcing $k = \Omega\bigl(\eps_0^{-2}\log N\bigr)$; see \cref{sec:whydeflation}.
\end{remark}

\section{Proof of the main theorem}\label{sec:mainproof}

\begin{proof}[Proof of \cref{thm:main}]
Take the pivot $\bar{\vec{A}} \coloneqq \bar{\vec{A}}_0$, so that $\vec{B} = \vec{A} - \bar{\vec{A}}_0$ and $\vec{B}_i = \vec{A}_i - \bar{\vec{A}}_0$.
Let $\eps_0$ be as in \cref{eq:eps0}, run \textsc{EstimateScores}$(\bar{\vec{A}}_0, \eps_0, \delta)$, and output $\hat{i} \in \argmin_i (\|\vec{B}_i\|_\F^2 - 2\hat{\varphi}_i)$.
\Cref{asm:warm} supplies $\|\vec{B}\|_\F \le \gamma\OPT$, so $\bar{\vec{A}}_0$ is a legal pivot for \cref{lem:select}.
By \cref{lem:estimate}, the hypothesis \cref{eq:relacc} of \cref{lem:select} holds with probability at least $1-\delta$, and on that event \cref{lem:select} gives $d_{\hat{i}} \le (1+\eps)\OPT$.
The cost is $O(k)$ queries, with
\[
k = O\bigl( \eps_0^{-2}\sqrt{\log(N/\delta)} + \log(1/\delta) \bigr)
= O\bigl( \gamma^4\eps^{-2}\sqrt{\log(N/\delta)} + \log(1/\delta) \bigr),
\]
since $\eps_0^{-2} = 4\gamma^2(1+\gamma)^2\eps^{-2} = O(\gamma^4\eps^{-2})$, using $1 + \gamma \le 2\gamma$.
\end{proof}

\section{Discussion}\label{sec:discussion}

\subsection{Left queries are necessary, in any bounded-precision model}\label{sec:onesided}

Fix $m$ and a code $\{\vec{v}_1, \dots, \vec{v}_N\} \subset \{\pm 1\}^m$ with pairwise Hamming distance at least $m/4$; by Gilbert--Varshamov such a code exists with $m = O(\log N)$.
Embed each $\vec{v}_i$ in $\mathbb{R}^n$ by zero-padding and take
\[
\cF \coloneqq \{\vec{e}_1\vec{v}_i^\T\}_{i \in [N]},
\qquad
\vec{A} \coloneqq \vec{e}_1\vec{w}^\T + \sqrt{m}\,\vec{e}_2\vec{e}_3^\T
\quad\text{for some codeword } \vec{w} .
\]
The second term of $\vec{A}$ lies outside the span of $\cF$ and is there only to keep $\OPT > 0$, as \cref{asm:warm} requires.
Since $\vec{A} - \vec{e}_1\vec{v}_i^\T = \vec{e}_1(\vec{w} - \vec{v}_i)^\T + \sqrt{m}\,\vec{e}_2\vec{e}_3^\T$ has orthogonal summands, and $\|\vec{w} - \vec{v}_i\|_2^2 = 4\operatorname{Ham}(\vec{w}, \vec{v}_i) \ge m$ for $\vec{v}_i \ne \vec{w}$,
\[
d_i^2 = \|\vec{w} - \vec{v}_i\|_2^2 + m,
\qquad
\OPT = \sqrt{m},
\qquad
d_i \ge \sqrt{2m} = \sqrt{2}\,\OPT \ \text{ for } \vec{v}_i \ne \vec{w}.
\]
A $(1+\eps)$-approximation with $\eps < \sqrt2 - 1$ therefore forces exact identification of $\vec{w}$.
The pivot $\bar{\vec{A}}_0 = \vec{0}$ is a legal warm start with $\gamma = \sqrt{2}$, since $\|\vec{A}\|_\F = \sqrt{2m}$, and it carries no information about $\vec{w}$.

A right query $\vec{y}$ returns $\vec{A}\vec{y} = (\vec{w}^\T\vec{y})\,\vec{e}_1 + \sqrt{m}\,y_3\,\vec{e}_2$, whose second component is known in advance.
If query entries are restricted to $\{\pm1, 0\}$ then $\vec{w}^\T\vec{y} \in \{-m, \dots, m\}$, so each right query has at most $2m+1$ possible outcomes.
An algorithm making $T$ right queries, each chosen as a function of all previous answers, is a decision tree of depth $T$ and branching factor $2m+1$, and must have at least $N$ leaves; hence
\[
T \ge \frac{\log N}{\log(2m+1)} = \Omega\!\left( \frac{\log N}{\log\log N} \right).
\]
The same bound holds if answers are merely observed to $O(\log m)$ bits of precision.
By contrast, a single left query with $\vec{x} = \vec{e}_1$ returns $\vec{A}^\T\vec{e}_1 = \vec{w}$ in full.

Thus one-sided access costs $\widetilde{\Omega}(\log N)$ on instances that two-sided access solves in one query, and this separation persists at the level of the general problem: \cref{thm:main} gives $O(\sqrt{\log N})$ two-sided queries for every instance.

Two caveats.
First, with unrestricted real entries and exact arithmetic the lower bound is void: the single right query $\vec{y} = (1, 2^{-1}, 2^{-2}, \dots)$ encodes all of $\vec{w}$ in one real number.
Lower bounds in this model are only meaningful under sign-restricted queries, finite precision, or additive noise; our upper bound, by contrast, survives all of these up to the modifications of \cref{sec:signs}.
Second, the instance shows that any algorithm achieving $o(\log N)$ queries must exploit two-sidedness on \emph{spiky} (low-stable-rank) differences $\vec{A} - \vec{A}_i$; this is exactly what drives the design of \cref{sec:estimate}.

\subsection{Why deflation: the variance obstruction}\label{sec:whydeflation}

The natural one-shot estimator for $\inner{\vec{B}}{\vec{B}_i}$ is a bilinear sketch, $\frac{1}{k^2}\inner{\vec{X}^\T\vec{B}_i\vec{Y}}{\vec{X}^\T\vec{B}\vec{Y}}$, or its one-sided cousin $\frac{1}{k}\inner{\vec{B}_i\vec{Y}}{\vec{B}\vec{Y}}$.
For the one-sided version, \cref{lem:hw} applies with $\vec{N}_i = \vec{B}_i^\T\vec{B}$ and the sub-Gaussian exponent is $k t^2 / \|\vec{B}_i^\T\vec{B}\|_\F^2$.
When $\vec{B}_i$ and $\vec{B}$ are both spiky and share a left factor --- say $\vec{B}_i = \vec{u}\vec{a}_i^\T$ and $\vec{B} = \vec{u}\vec{b}^\T$ with $\|\vec{u}\|_2 = 1$, which is the situation in \cref{sec:onesided} with $\vec{u} = \vec{e}_1$ --- one has $\|\vec{B}_i^\T\vec{B}\|_\F = \|\vec{B}_i\|_\F\|\vec{B}\|_\F$.
At the target accuracy $t = \eps_0\|\vec{B}_i\|_\F\|\vec{B}\|_\F$ the exponent then degrades to $\eps_0^2 k$, and the union bound over $N$ hypotheses forces $k = \Omega\bigl(\eps_0^{-2}\log N\bigr)$.
Equivalently: the effective sample size of a bilinear sketch of a rank-one matrix is $k$, not $k^2$, because $\|\vec{X}^\T\vec{u}\vec{a}^\T\vec{Y}\|_\F^2 = \|\vec{X}^\T\vec{u}\|_2^2\,\|\vec{Y}^\T\vec{a}\|_2^2$ is a product of two $k$-degree-of-freedom factors.

The routine of \cref{sec:estimate} removes both spiky factors before sketching.
The candidate factor is flattened offline, by truncating the known matrix $\vec{B}_i$ at rank $r_0$; the truncation error is controlled \emph{deterministically} against $\|\vec{R}\|_2$ through the nuclear norm, in \cref{eq:lowrank}.
The unknown factor is flattened by the rangefinder.
It is worth being precise about where left queries enter here.
The basis $\vec{Q} = \orth(\vec{B}\vec{\Omega})$ is computable from right queries alone, and \cref{eq:deflation} is a statement about $\vec{Q}$ only; on the rank-one instances of \cref{sec:onesided} one even has $(\vec{I} - \vec{Q}\vec{Q}^\T)\vec{B} = \vec 0$ exactly.
What one-sided access cannot deliver is the deflated matrix $\widehat{\vec{B}} = \vec{Q}\vec{Q}^\T\vec{B}$ itself: forming $\vec{Q}^\T\vec{B}$ is precisely a left-query operation.
Without $\widehat{\vec{B}}$ the algorithm can neither evaluate the leading term $\inner{\vec{B}_i}{\widehat{\vec{B}}}$ of $\hat\varphi_i$ nor form the residual sketch $\vec{R}\vec{Y}$, and the estimator collapses back to the un-deflated one.

After both flattening steps, both Hanson--Wright branches carry exponents of order $\eps_0^4 k^2$ and $\eps_0^2 k^2$ (\cref{rem:k2}), and the union bound over $N$ hypotheses is absorbed at $k = \Theta\bigl(\eps_0^{-2}\sqrt{\log N}\bigr)$.

\subsection{Sign-restricted queries}\label{sec:signs}

Suppose query entries must lie in $\{\pm 1, 0\}$, and suppose an a priori bound $\|\vec{A}\|_2 \le \Lambda$ is known.
The algorithm survives with two modifications, at the cost of adjusting the absolute constants in \cref{eq:kdef}.

First, the Gaussian sketches $\vec{\Omega}$ and $\vec{Y}$ are replaced by Rademacher sketches.
\Cref{lem:hw} holds verbatim for Rademacher vectors \cite{rudelson2013hanson}.
\Cref{lem:rsvd} is a Gaussian statement, and we are not aware of a reference establishing \cref{eq:deflation} verbatim for Rademacher $\vec{\Omega}$; we take the sub-Gaussian analogue as a hypothesis, and flag it as a technical gap.

Second, the left queries with the real-valued columns of $\vec{Q}$ are simulated by binary expansion.
Writing $\vec{q}_j$ to $b$ bits as a signed combination $\tilde{\vec{q}}_j = \sum_{a=1}^{b} 2^{-a}\vec{s}_{j,a}$ with $\vec{s}_{j,a} \in \{\pm1,0\}^n$ costs $b$ legal left queries per column and yields $\|\vec{q}_j - \tilde{\vec{q}}_j\|_2 \le 2^{-b}\sqrt{n}$.
The algorithm then forms $\widetilde{\vec{W}}$ with rows $(\vec{A}^\T\tilde{\vec{q}}_j)^\T - \vec{q}_j^\T\bar{\vec{A}}$ and sets $\widehat{\vec{B}} \coloneqq \vec{Q}\widetilde{\vec{W}}$, so that
\[
\vec{E} \coloneqq \widehat{\vec{B}} - \vec{Q}\vec{Q}^\T\vec{B},
\qquad
\|\vec{E}\|_\F = \|\widetilde{\vec{W}} - \vec{W}\|_\F \le 2^{-b}\sqrt{nk}\,\Lambda .
\]
Now $\vec{R} \coloneqq \vec{B} - \widehat{\vec{B}}$ no longer satisfies $\vec{R} = (\vec{I} - \vec{Q}\vec{Q}^\T)\vec{R}$.
The splitting identity $\inner{\vec{B}_i}{\vec{R}} = \inner{\vec{B}_i^L}{\vec{R}} + \inner{\vec{B}_i^H}{\vec{R}}$ is unaffected, since $\vec{B}_i^L + \vec{B}_i^H = \vec{B}_i$ holds by construction; the perturbation instead enters the leading term of $\hat\varphi_i$, which now estimates $\inner{\vec{B}_i}{\vec{Q}\vec{Q}^\T\vec{B}}$ with error $\inner{\vec{B}_i}{\vec{E}}$, of magnitude at most $\|\vec{B}_i\|_\F \|\vec{E}\|_\F$.
Likewise $\|\vec{R}\|_\F \le \|\vec{B}\|_\F + \|\vec{E}\|_\F$ and, on the deflation event, $\|\vec{R}\|_2 = O\bigl(\|\vec{B}\|_\F/\sqrt{k}\bigr) + \|\vec{E}\|_\F$.
All three perturbations are absorbed once $\|\vec{E}\|_\F = O(\eps_0\|\vec{B}\|_\F)$ with a small enough implied constant.

Setting $b$ requires a lower bound on $\|\vec{B}\|_\F$, which is obtained from the right queries already made: $\widehat{\beta}^2 \coloneqq \frac{1}{k}\|\vec{B}\vec{\Omega}\|_\F^2$ is an averaged quadratic form with matrix $\vec{B}^\T\vec{B}$, mean $\|\vec{B}\|_\F^2$, $\|\vec{B}^\T\vec{B}\|_\F \le \|\vec{B}\|_2\|\vec{B}\|_\F \le \|\vec{B}\|_\F^2$, and $\|\vec{B}^\T\vec{B}\|_2 \le \|\vec{B}\|_\F^2$, so \cref{lem:hw} with $t = \|\vec{B}\|_\F^2/2$ gives $\tfrac12\|\vec{B}\|_\F^2 \le \widehat{\beta}^2 \le 2\|\vec{B}\|_\F^2$ except with probability $e^{-\Omega(k)}$.
(If $\widehat{\beta} = 0$ then $\vec{B}\vec{\Omega} = \vec{0}$, which for $\|\vec{B}\|_\F > 0$ has probability zero; the algorithm outputs $\argmin_i \|\vec{B}_i\|_\F$.)
Taking $b = O\bigl( \log(nk\Lambda/(\eps_0\|\vec{B}\|_\F)) \bigr)$ suffices.

The net effect is a multiplicative $O(\log(n\Lambda/(\eps\,\OPT)))$ overhead on the left queries, which are one third of the total; the right-query count is unchanged.

\subsection{Computational cost}\label{sec:computation}

The query count is the complexity measure here, but the offline work is nontrivial: the algorithm computes a top-$r_0$ SVD of $\vec{B}_i$ for every $i \in [N]$.
Since $r_0 \le k$, randomized low-rank SVD brings the per-candidate cost to $O(n^2 k)$, for $O(N n^2 k)$ in total.
No attempt has been made to optimize this.

\subsection{Open problems}\label{sec:open}

\begin{enumerate}
\item \textbf{Is the warm start free?}
\Cref{asm:warm} is the only nonstandard input.
A Scheff\'e-style tournament at constant accuracy would produce a constant-factor pivot, but it requires each pairwise comparison to succeed with confidence $1 - N^{-2}$, and any estimator built from $k$ scalar samples of an \emph{un-deflated} residual concentrates only at rate $e^{-\Theta(k)}$ on spiky instances, forcing $k = \Omega(\log N)$.
The deflation that repairs the rate is pivot-relative: its error scale is $\eps_0\|\vec{B}_i\|_\F\|\vec{B}\|_\F$, and $\|\vec{B}\|_\F$ is the \emph{current} pivot's distance, not $\OPT$.
Iterated contraction from an arbitrary $\vec{A}_1$ therefore costs $O\bigl(\sqrt{\log N}\,\log(2+\kappa)\bigr)$ queries with $\kappa = \|\vec{A} - \vec{A}_1\|_\F/\OPT$.
Can the $\log\kappa$ be removed in a finite-precision model, giving a scale-free $O(\eps^{-2}\sqrt{\log N})$ algorithm?
\item \textbf{Matching lower bound.}
We know of no $\Omega(\sqrt{\log N})$ lower bound for two-sided sign-restricted (or finite-precision) queries.
The instances of \cref{sec:onesided} are solved in one two-sided query, and Fano-type arguments leak because each query returns $n$ coordinates.
\item \textbf{Dependence on $\eps$.}
The error term in \cref{lem:estimate} is already localized as $\eps_0\|\vec{B}_i\|_\F\|\vec{B}\|_\F$, which suggests the $\eps^{-2}$ might improve to $\eps^{-1}$ by a sharper final round; we do not currently see how.
\item \textbf{One-sided complexity.}
For right-only queries the truth plausibly is $\Theta(\eps^{-2}\log N)$; the upper bound follows from the un-deflated estimator of \cref{sec:whydeflation}, while \cref{sec:onesided} gives $\Omega(\log N/\log\log N)$.
\end{enumerate}

\printbibliography

\end{document}

